%% exemplo-relatorio.tex
%% Exemplo de uso da classe uftreport
%%
%% Compilar com:
%%   pdflatex exemplo-relatorio.tex
%%   bibtex   exemplo-relatorio   (se houver referências BibTeX)
%%   pdflatex exemplo-relatorio.tex
%%   pdflatex exemplo-relatorio.tex
%%
%% São necessárias ao menos duas passagens do pdflatex para que
%% a paginação (via zref) e as referências cruzadas fiquem corretas.
%%
%% ── Opções de tipo de documento ───────────────────────────────────────
%%   report           → trabalho de disciplina ou relatório (padrão)
%%   projectresearch  → projeto de pesquisa (PIBIC, TCC, etc.)
%%   projectextension → projeto de extensão
%%   manual           → manual técnico (cabeçalhos de capítulo maiores)
%%   techreport       → relatório técnico numerado (exibe número TR na capa)
%%
%% ── Cursos disponíveis em \department{} ───────────────────────────────
%%   CC  → Ciência da Computação (padrão)
%%   LC  → Licenciatura em Computação
%%   EC  → Engenharia Civil
%%   EE  → Engenharia Elétrica
%%
%% ── Observações sobre o abntex2cite ───────────────────────────────────
%%   Não inclua \renewcommand{\backrefpagesname}, \renewcommand{\backref}
%%   nem \renewcommand*{\backrefalt} — a classe não carrega o pacote
%%   backref e esses comandos não existem.

\documentclass[techreport]{uftreport}

\usepackage[alf,abnt-emphasize=bf]{abntex2cite}

% ── Metadados do documento ─────────────────────────────────────────────
\title{Título do Trabalho ou Relatório}
\foreigntitle{Title of the Work or Report}

\author{Nome}{Sobrenome}
% \author{Segundo}{Autor}   % repita para cada coautor

\advisor{Prof.}{João}{Silva}{Dr.}
% \advisor{Profa.}{Maria}{Souza}{Dra.}   % repita para coorientador(a)

\department{CC}   % CC | LC | EC | EE

\class{Algoritmos e Estruturas de Dados II -- 2024.1}
% Remova \class{} se o documento não for vinculado a uma disciplina

% Número TR — usado apenas com a opção techreport
\TRNumber{001}

\date{1}{6}{2026}   % dia, mês, ano

% ── Palavras-chave (aparecem ao final do Resumo e do Abstract) ─────────
\keyword{Computação}
\keyword{Algoritmos}
\keyword{LaTeX}

\foreignkeyword{Computer Science}
\foreignkeyword{Algorithms}
\foreignkeyword{LaTeX}

% ======================================================================
\begin{document}

% ── Elementos pré-textuais ─────────────────────────────────────────────
\frontmatter        % desativa numeração de páginas, estilo vazio

\maketitle          % capa no estilo UFT/CC

\makefrontpage      % folha de rosto com texto gerado pela classe
\newpage

% \dedication{À minha família, pelo apoio incondicional.}

% \begin{acknowledgement}
%   Agradeço ao professor orientador pela paciência e dedicação.
% \end{acknowledgement}

\begin{abstract}
Este documento exemplifica o uso da classe \texttt{uftreport} para
a produção de relatórios técnicos, manuais e trabalhos de disciplinas
do Curso de Ciência da Computação da Universidade Federal do Tocantins.
A classe oferece cinco modos de uso --- \texttt{report},
\texttt{projectresearch}, \texttt{projectextension}, \texttt{manual}
e \texttt{techreport} --- e incorpora o padrão visual institucional da UFT.
\end{abstract}

\begin{foreignabstract}
This document illustrates the use of the \texttt{uftreport} class for
producing technical reports, manuals and course assignments for the
Computer Science program at the Universidade Federal do Tocantins.
\end{foreignabstract}

\tableofcontents
\listoffigures    % remova se não houver figuras
\listoftables     % remova se não houver tabelas

% ── Texto principal ────────────────────────────────────────────────────
\mainmatter         % ativa numeração árabe

% \ChapterStart marca a página em que a numeração deve começar.
% Deve ser chamado imediatamente antes do primeiro \chapter.
\ChapterStart{first}{Introdução}

% ======================================================================
\chapter{Introdução}

Este é o primeiro capítulo. O texto segue as normas de formatação
definidas na classe, com espaçamento 1,5, margens ABNT e cabeçalhos
no estilo UFT/CC.

\section{Contexto}

Texto da seção. A seção é formatada com \textbf{negrito azul} e uma
linha horizontal discreta abaixo do título.

\subsection{Motivação}

Texto da subseção.

\subsubsection{Objetivo específico}

Texto de subsubseção.

% ======================================================================
\chapter{Fundamentação Teórica}

\section{Exemplo de código-fonte}

O ambiente \texttt{lstlisting} está pré-configurado com destaque
sintático e numeração de linhas:

\begin{lstlisting}[language=Python, caption={Busca binária em Python}]
def busca_binaria(vetor, alvo):
    esq, dir = 0, len(vetor) - 1
    while esq <= dir:
        meio = (esq + dir) // 2
        if vetor[meio] == alvo:
            return meio
        elif vetor[meio] < alvo:
            esq = meio + 1
        else:
            dir = meio - 1
    return -1
\end{lstlisting}

\section{Exemplo de tabela}

\begin{table}[ht]
\centering
\caption{Comparação de algoritmos de ordenação}
\begin{tabular}{lccl}
\toprule
\textbf{Algoritmo} & \textbf{Melhor} & \textbf{Pior} & \textbf{Estável} \\
\midrule
Bubble Sort   & $O(n)$       & $O(n^2)$      & Sim \\
Merge Sort    & $O(n\log n)$ & $O(n\log n)$  & Sim \\
Quick Sort    & $O(n\log n)$ & $O(n^2)$      & Não \\
Heap Sort     & $O(n\log n)$ & $O(n\log n)$  & Não \\
\bottomrule
\end{tabular}
\label{tab:algoritmos}
\end{table}

% ======================================================================
\chapter{Metodologia}

Descreva aqui a metodologia utilizada.

% ======================================================================
\chapter{Resultados}

Apresente os resultados obtidos.

% ======================================================================
\chapter{Conclusão}

Sintetize as conclusões do trabalho \citeonline{lamport1994}~\cite{cormen2022}.

% ── Elementos pós-textuais ─────────────────────────────────────────────
\backmatter

% Referências com BibTeX + abntex2cite:
\bibliography{references}

% Alternativa sem BibTeX (referências manuais):
% \begin{thebibliography}{9}
%   \bibitem{cormen2022}
%     CORMEN, T. H. et al.
%     \textit{Introdução a Algoritmos}. 4. ed. MIT Press, 2022.
%   \bibitem{lamport1994}
%     LAMPORT, L.
%     \textit{\LaTeX: A Document Preparation System}. 2. ed.
%     Addison-Wesley, 1994.
% \end{thebibliography}

\end{document}